Per a mesurar la velocitat del so en l'aire podem fer servir un tub de ressonància. Regulant el nivell de l'aigua, es poden produir situacions de ressonància quan l'ona estacionària té un ventre a l'extrem obert del tub. Quan el diapasó vibra amb una freqüència de 440~Hz, fem baixar el nivell de l'aigua fins que observem la primera situació de ressonància per a $h = 19\ \mathrm{cm}$, que es reconeix perquè es produeix una intensificació nítida del so, i també observem una segona situació de ressonància per a $h = 57\ \mathrm{cm}$.

\figura[0.22]{fig1.png}

\begin{apartat}{1}
Dibuixeu l'esquema de l'ona estacionària per a cadascuna de les situacions de ressonància descrites i determineu la velocitat del so en l'aire.
\end{apartat}

\begin{apartat}{1}
Si el diapasó emet ones sonores amb una potència de 0,01~W, calculeu els decibels que percebrà una persona situada a 3~m.
\end{apartat}

\blocetiquetat{Dada:}{Intensitat del llindar d'audició: $I_0 = 10^{-12}\ \mathrm{W\,m^{-2}}$}
